%! Solving by Square Roots \& Complex Numbers — Unit 2, Lesson 2.3 — Student Packet Cover Sheet
% SLIDES-FIRST shape (2026-09-25), condensed to three rounds (2026-09-29). 12pt; \coverbanner
% measures the title block. TWO scored rows: the slide handout (the student's notes, the pages
% right after this cover) and the homework — a DeltaMath set assigned ONLINE (user choice), so it
% is not in this packet. Homework is due the first class after two study halls — never "next class".
\documentclass[12pt]{article}
\usepackage{algebra2-article}
\usepackage{algebra2-boxes}

\begin{document}

\coverbanner{Unit 2 \quad Quadratic Functions}{Lesson 2.3 \quad Solving by Square Roots \& Complex Numbers}

\namedateperiod
\vspace{0.18in}

\boxguard[14]
\begin{learningtargetbox}
\small
\begin{enumerate}[leftmargin=1.5em, itemsep=5pt]
  \item \ldots{} solve $x^2=k$ and $(x-h)^2=k$ with the \textbf{Square Root Property}, keeping
        both roots and writing them in \textbf{simplest radical form}.
  \item \ldots{} use the \textbf{imaginary unit} $i=\sqrt{-1}$ to rewrite $\sqrt{-k}$ and to write
        a \textbf{complex number} in \textbf{standard form} $a+bi$, naming its \textbf{real part}
        and \textbf{imaginary part}.
  \item \ldots{} add, subtract, and multiply complex numbers, replacing $i^2$ with $-1$, and
        multiply \textbf{complex conjugates} to get a real number.
\end{enumerate}
\end{learningtargetbox}

\vspace{0.14in}

\boxguard[14]
\begin{tocbox}
\small
\renewcommand{\arraystretch}{1.5}
\begin{tabularx}{\linewidth}{@{} c l X r @{}}
\rowcolor{forestbg}
\textbf{\#} & \textbf{Component} & \textbf{Description} & \textbf{Score} \\
\midrule
1 & Slide Notes & The printed slides: the warm-up, three rounds of definitions, examples, and
                  practice ladders (square roots, the imaginary unit, complex arithmetic), and
                  the Final Round --- worked in the notes column & \blank{1.2cm} \\
2 & DeltaMath   & \textbf{2.3 Square Roots \& Complex Numbers} --- online, not in this packet;
                  \textbf{scored; due the first class after two study halls} & \blank{1.2cm} \\
\midrule
  & \multicolumn{2}{r}{\textbf{Total}} & \blank{1.2cm} \\
\end{tabularx}
\end{tocbox}

\vspace{0.14in}

\boxguard[8]
\begin{remindbox}
\small
The \textbf{Square Root Property}: if $x^2=k$ and $k\ge0$, then $x=\pm\sqrt{k}$ --- isolate the
square first, keep \emph{both} roots, and simplify ($\sqrt{12}=2\sqrt3$). The \textbf{imaginary
unit} is $i=\sqrt{-1}$, so $i^2=-1$ and $\sqrt{-k}=i\sqrt{k}$. A \textbf{complex number} $a+bi$
has real part $a$ and imaginary part $b$. Add or subtract by combining like parts; multiply by
distributing, then replace $i^2$ with $-1$. \textbf{Conjugates} give a real product:
$(a+bi)(a-bi)=a^2+b^2$. Because $i^2=-1$, $\sqrt{-a}\cdot\sqrt{-b}$ is \emph{not} $\sqrt{ab}$ ---
pull the $i$ out of every negative root \emph{before} you multiply.
\end{remindbox}

\end{document}
